% ===========================================================================
% exp15_som.tex - complete code listings for Experiment 15
% Source: notebooks/07_self_organizing_map.ipynb (executed)
% Repository: https://github.com/Atik203/ML-Algorithms
% ===========================================================================

\begin{codebox}[title={Core numerical and plotting libraries}]
# ---- Core numerical and plotting libraries ----
import numpy as np
import pandas as pd
import matplotlib.pyplot as plt
import seaborn as sns

# ---- Dataset + scaling ----
from sklearn.datasets import load_wine
from sklearn.preprocessing import MinMaxScaler
# ---- MiniSom: lightweight Kohonen Self-Organizing Map implementation ----
from minisom import MiniSom

# Styling setup (consistent look across all notebooks)
plt.style.use('seaborn-v0_8-whitegrid' if 'seaborn-v0_8-whitegrid' in plt.style.available else 'default')
plt.rcParams['font.sans-serif'] = 'DejaVu Sans'
plt.rcParams['figure.figsize'] = (10, 6)
plt.rcParams['figure.dpi'] = 110

np.random.seed(42)
print("MiniSom and libraries loaded successfully!")

# ---- Load the Wine recognition dataset ----
wine = load_wine()
X_raw = wine.data          # 178 samples x 13 chemical features
y = wine.target            # 3 wine cultivars (only used for visualization - SOM is unsupervised)
feature_names = wine.feature_names
target_names = wine.target_names

print(f"Samples: {X_raw.shape[0]}, Features: {X_raw.shape[1]}, Classes: {target_names.tolist()}")

# ---- MinMax scaling to [0, 1] ----
# SOM weight vectors live in the same range as the inputs, so scaling keeps learning stable.
scaler = MinMaxScaler()
X = scaler.fit_transform(X_raw)
\end{codebox}

\begin{codebox}[title={Configure the Kohonen grid}]
# ---- Configure the Kohonen grid ----
grid_x, grid_y = 12, 12  # 12x12 grid = 144 neurons (fewer than the 178 samples)
input_len = X.shape[1]   # 13 chemical features per neuron weight vector

som = MiniSom(
    x=grid_x,
    y=grid_y,
    input_len=input_len,
    sigma=1.5,                       # initial neighborhood radius
    learning_rate=0.5,               # initial learning rate
    neighborhood_function='gaussian',# smooth Gaussian neighborhood kernel
    random_seed=42
)

# ---- Initialize weights with PCA for topologically faithful, faster convergence ----
som.pca_weights_init(X)
print("Initial Quantization Error:", som.quantization_error(X))

# ---- Train: 5000 random samples presented to the map ----
som.train_random(data=X, num_iteration=5000, verbose=False)
final_qe = som.quantization_error(X)   # average distance from samples to their BMU
final_te = som.topographic_error(X)    # fraction of samples whose 2nd BMU is not a grid neighbor

print(f"Training complete after 5,000 iterations!")
print(f" - Final Quantization Error (mean distance to BMU): {final_qe:.4f}")
print(f" - Final Topographic Error (proportion of foldings): {final_te:.4f}")
\end{codebox}

\begin{codebox}[title={Plot the U-Matrix (unified distance map)}]
# ---- Plot the U-Matrix (unified distance map) ----
# Each cell = average distance between a neuron and its immediate grid neighbors.
# Dark valleys = dense clusters; bright ridges = boundaries between clusters.
plt.figure(figsize=(10, 8))
plt.pcolor(som.distance_map().T, cmap='bone_r', alpha=0.9)
cbar = plt.colorbar()
cbar.set_label("Normalized Inter-Neuron Distance (U-Matrix)", fontsize=11)

# ---- Project every sample onto its Best Matching Unit (BMU) ----
markers = ['o', 's', '^']
colors = ['red', 'green', 'blue']

for idx, x_vec in enumerate(X):
    w = som.winner(x_vec)  # grid coordinates of the best matching neuron for this sample
    plt.plot(
        w[0] + 0.5,        # +0.5 centers the marker inside the pcolor cell
        w[1] + 0.5,
        markers[y[idx]],   # marker shape encodes the true cultivar
        markerfacecolor='None',
        markeredgecolor=colors[y[idx]],
        markersize=10,
        markeredgewidth=2
    )

# Build a legend for the three cultivars
for c_idx, c_name in enumerate(target_names):
    plt.plot([], [], marker=markers[c_idx], color=colors[c_idx], linestyle='None',
             label=f'Cultivar {c_name}', markersize=9, markeredgewidth=2)

plt.title("SOM U-Matrix with Projected Wine Cultivars", fontsize=14, fontweight='bold')
plt.legend(bbox_to_anchor=(1.25, 1), loc='upper left')
plt.tight_layout()
plt.show()
\end{codebox}
